\predpoglavjevkazalu{Načrt ravnanja z raziskovalnimi podatki (NRRP)}\label{cha:nrrp}
Ta priloga je obvezna in jo mora kandidat izpolniti ob predstavitvi rezultatov raziskovalnega dela oziroma ob oddaji doktorske disertacije (ta stavek pred oddajo izbrišite).

\begingroup
\setlist[enumerate]{label=\arabic*.,leftmargin=7.5mm,topsep=2pt,itemsep=12pt,parsep=0pt}
% Vsako polje ima na dnu prostor za odgovor.
\newcommand{\nrrpPolje}[1]{\vspace*{4pt}#1\par\vspace*{12pt}}
\setlength{\tabcolsep}{6pt}
\noindent\begin{tabular}{|>{\setlength{\parskip}{8pt}}p{\dimexpr\textwidth-2\tabcolsep-2\arrayrulewidth\relax}|}
\hline
\nrrpPolje{%
  \textbf{Ime in priimek doktoranda/doktorandke:} \avtor

  \textbf{Doktorski program in področje:} \doktorskiProgram

  \textbf{Naslov doktorske disertacije:} \naslovSlo}\\
\hline
\nrrpPolje{%
  \textbf{Tip podatkov in metode njihovega zbiranja in/ali ustvarjanja}
  \begin{enumerate}
  \item Katere podatke ste zbirali oziroma pridobivali in/ali ustvarjali?
  \item Na kakšen način ste zbirali oziroma pridobivali in/ali ustvarjali nove podatke in kako ste uporabljali že obstoječe za potrebe vaše doktorske disertacije?
  \item Ali ste delali z občutljivimi podatki? Če da, kako ste poskrbeli za etično pridobivanje in/ali ustvarjanje podatkov?
  \end{enumerate}}\\
\hline
\nrrpPolje{%
  \textbf{Način hranjenja in zaščita podatkov med raziskovalnim delom za doktorsko disertacijo}
  \begin{enumerate}
  \item Na kakšen način ste hranili podatke?
  \item Če ste delali z občutljivimi podatki, kako ste skrbeli za njihovo varovanje in zaščito? (v nasprotnem primeru to vprašanje preskočite)
  \end{enumerate}}\\
\hline
\nrrpPolje{%
  \textbf{Dolgotrajna dostopnost in hranjenje podatkov}
  \begin{enumerate}
  \item Kje oziroma v katerem podatkovnem repozitoriju boste hranili podatke na dolgi rok po zaključku raziskovalnega dela in omogočili njihovo dostopnost v skladu z zahtevo 50.~člena Pravilnika o doktorskem študiju UL?
  \item Ali načrtujete za določen čas omejiti dostop do podatkov? Če da, pojasnite razloge (npr.~zaradi zaščite intelektualne lastnine ali patenta, ali navedite drug razlog).
  \end{enumerate}}\\
\hline
\end{tabular}
\endgroup
